% =====================================================================
%  Portada y consideraciones generales de uso de la plantilla
% =====================================================================

\thispagestyle{uvestilo}

\begin{instruccion}[Consideraciones generales para el uso de esta plantilla]
\begin{itemize}
  \item Cuide la redacción y la ortografía; respete tipografía, interlineado y justificación.
  \item No pierda el encabezado ni el pie de página al agregar contenido.
  \item Actualice los cuadros de \emph{Historial} y \emph{Revisiones} en cada entrega.
  \item El índice, los índices de cuadros y figuras y las referencias cruzadas se generan
        solos: basta con recompilar el documento dos veces.
  \item \textbf{Todos} los cuadros y figuras deben estar referidos desde el texto
        (\texttt{\textbackslash ref\{...\}}), indicando qué información presentan y una
        breve explicación.
  \item Para producir la versión final de entrega basta con ejecutar \texttt{make final}:
        las guías verdes y los ejemplos azules desaparecen y el resto del documento no
        cambia. Para hacerlo permanente, ponga a \texttt{false} los dos interruptores
        \texttt{mostrarinstrucciones} y \texttt{mostrarejemplos} de
        \texttt{config/metadata.tex}.
  \item Sustituya cada marcador \campo{así} por el contenido que corresponda.
\end{itemize}
\end{instruccion}

\begin{notanormativa}[Cómo leer los bloques de color de esta plantilla]
\begin{itemize}
  \item \textcolor{instrverde}{\textbf{Verde}} — guía de la plantilla: explica qué debe
        escribirse. Se elimina en la versión final.
  \item \textcolor{ejemazul}{\textbf{Azul}} — ejemplo ilustrativo. No es contenido a conservar.
  \item \textcolor{polinaranja}{\textbf{Naranja}} — política académica del curso.
        Es una regla \emph{local}, no un requisito de \norma.
  \item \textcolor{normgris}{\textbf{Gris}} — referencia normativa o aclaración conceptual,
        con su fuente identificada.
\end{itemize}
Esta separación responde a una observación explícita de la revisión: el lector debe poder
distinguir de inmediato un requisito del curso de una recomendación de la norma.
\end{notanormativa}

\clearpage

% --------------------------- PORTADA ---------------------------------
\begin{center}
  \vspace*{1.5\baselineskip}

  {\large\bfseries \uvUniversidad}\\[2pt]
  {\normalsize \uvFacultad}\\[2pt]
  {\normalsize \uvPrograma}\\[2pt]
  {\normalsize \uvEE}

  \vspace{3\baselineskip}

  {\Huge\bfseries\color{uvazul} \docTitulo}

  \vspace{1.2\baselineskip}
  {\LARGE \docProyecto}

  \vspace{3\baselineskip}

  \begin{tabularx}{0.82\textwidth}{|E{3.6cm}|Y|}
    \hline
    Fecha            & \docFecha \\ \hline
    Equipo           & \docEquipo \\ \hline
    Integrantes      & \docIntegrantes \\ \hline
    Versión del documento & \docVersionDoc \\ \hline
  \end{tabularx}

  \vspace{2.5\baselineskip}

  \begin{minipage}{0.86\textwidth}\small\centering
    El presente documento se basa en la estructura de documentación de la prueba indicada en
    \normap{3}, \emph{Software and systems engineering — Software testing — Part~3: Test
    documentation}, y en el modelo de procesos de \normap{2}.
  \end{minipage}

  \vfill
\end{center}

\clearpage

% --------------------- INFORMACIÓN DEL PROYECTO ----------------------
\section*{Información del proyecto}
\phantomsection
\addcontentsline{toc}{part}{Información del proyecto}

\begin{instruccion}
Identifique el proyecto que será sometido a prueba y su situación actual. El
\emph{estatus} sirve para justificar decisiones posteriores: no se planifica igual la prueba
de un sistema terminado que la de uno en desarrollo activo.
\end{instruccion}

\begin{table}[H]
\centering
\caption{Identificación del proyecto sometido a prueba.}
\label{tab:info-proyecto}
\begin{tabularx}{\textwidth}{|E{5cm}|Y|}
  \hline
  Equipo desarrollador            & \campo{Quién construyó el software} \\ \hline
  Proyecto                        & \campo{Nombre y versión del sistema} \\ \hline
  Estatus actual del proyecto     & \campo{En desarrollo / entregado / en mantenimiento} \\ \hline
  Experiencia Educativa           & \uvEE \\ \hline
  Equipo de prueba                & \docEquipo \\ \hline
  Periodo del esfuerzo de prueba  & \campo{dd/mm/aaaa – dd/mm/aaaa} \\ \hline
\end{tabularx}
\end{table}

El Cuadro~\ref{tab:info-proyecto} identifica el sistema que será probado, quién lo desarrolló
y en qué estado se encuentra al iniciar el esfuerzo de prueba; estos datos delimitan lo que el
resto del plan puede asumir como disponible.
